\documentclass{in1150-innlevering}
\title{Innleveringsoppgave 4 i IN1150}
\author{Eilif Akerjordet}
\date{12. februar 2021}
\begin{document}
\maketitle

\section*{Oppgave 1}
  $S=\{1,2,3\}$
  $T=\{a,b,c,d\}$
\begin{deloppgaver}
  \item
    $g=\{\tuple{1,a},\tuple{2,c}\}$
    Dette er ikke en funksjon.
  \item
    $h=\{\tuple{1,d},\tuple{2,d},\tuple{3,d}\}$ Dette er en funksjon, men den er verken injektiv eller surjektiv.
  \item
    $i=\{\tuple{2,b},\tuple{1,a},\tuple{3,c}\}$ Dette er en injektiv funksjon.
  \item
    $j=\{\tuple{3,c},\tuple{2,a},\tuple{3,b},\tuple{1,d}\}$ Dette er ikke en funksjon.
\end{deloppgaver}

\section*{Oppgave 2}
\begin{deloppgaver}
\item $h=\{\tuple{1,c},\tuple{2,a},\tuple{3,a}\}$
\item $h[A]=\{c, a, a\}$
\end{deloppgaver}

\section*{Oppgave 3}
\begin{deloppgaver}
\item 
  Definisjonen av den operasjon er en funksjon som tar en verdi fra definisjonsområrdet, og returnerer en verdi fra samme mengde. Da $A$ er en delmengde i $B$, og $f$ er en funksjon fra $A$ til $B$, gjør dette $f$ til en operasjon.
\item ---
\item ---
\end{deloppgaver}

\section*{Oppgave 4}
\begin{deloppgaver}
\item ---
\item ---
\item ---
\end{deloppgaver}
\section*{Oppgave 5}
\begin{deloppgaver}
  \item $P(\{2,a\}) = \{\{2\},\{a\},\{2,a\},\{\emptyset\}\}$
  \item $P(\{1\}\cup\{\emptyset, 2\}) = \{\{1\},\{\emptyset\},\{2\},\{\emptyset, 2\}\}$
  \item $P(\{\emptyset, 4\}) = \{\{\emptyset\},\{4\},\{\emptyset, 4\}\}$
  \item $P(P(\{57\})) = \{57\}$
\end{deloppgaver}
\section*{Oppgave 6}
\begin{deloppgaver}
  \item 
    \begin{itemize}
      \item
          $\overline{A\cup B}$ tilsvarer alle elementer som ikke er i $A$ eller $B$.
      \item
        $\overline{A}$ er alle elementer som er i $B$ og alt annet som ikke er i $A$.
      \item
        $\overline{B}$ er alle elementer som er i $A$ og alt annet som ikke er i $B$.
      \item
        $\overline{A} \cap \overline{B}$ er derfor alle elementer som ikke er i A eller B.
      \item
        Dette medfører at et tilfeldig element valgt fra $\overline{A\cup B}$ vil også være et element i $\overline{A} \cap \overline{B}$ 
    \end{itemize}
  \item
    Hvis vi går baklengs gjennom forklaringen i oppgave a, ser vi hvorfor et element fra  $\overline{A} \cap \overline{B}$ også vil være et element i $\overline{A\cup B}$
  \item
    I forklaringene over er det irrelevant hva som ligger i A og B. Beviset holder derfor sant for alle mengder A og B.
\end{deloppgaver}

\section*{Oppgave 7}
---

\section*{Oppgave 8}
  \begin{deloppgaver}
    \item ---
    \item ---
    \item ---
  \end{deloppgaver}

\end{document}
